%---------------------------------------------------
% !TEX program = xelatex
% کلاس سند: article برای اسناد متنی، با اندازه کاغذ A4 و فونت پایه 16
\documentclass[a4paper,16pt]{article}

% بسته‌های ریاضی برای معادلات و نمادهای پیشرفته
\usepackage{amsmath, amssymb, amstext,siunitx}

% بسته رنگ برای استفاده از رنگ در متن یا محیط‌ها
\usepackage{xcolor}

% بسته گرافیک برای درج تصاویر
\usepackage{graphicx}

% تنظیمات حاشیه صفحه
\usepackage{geometry}

% بسته لینک‌دهی (هایپرلینک‌ها)
\usepackage{hyperref}

% بسته صفحه‌آرایی برای تنظیم هدر و فوتر
\usepackage{fancyhdr}

% بسته xepersian برای پشتیبانی از زبان فارسی با XeLaTeX
\usepackage{xepersian}

% تعیین فونت فارسی
\settextfont{B Nazanin}

% تنظیم اندازه حاشیه‌ها از هر طرف برابر با 1 اینچ
\geometry{margin=1in}

% تنظیم سبک fancy برای سربرگ و پاورقی
\pagestyle{fancy}

% سمت چپ هدر: مشخص‌کننده سری تمرین و نام درس
\lhead{تمرین سری دو ـ ریاضی فیزیک 1}

% وسط هدر: نام دانشکده
\chead{دانشکده فیزیک دانشگاه صنعتی خواجه نصیرالدین طوسی}

% سمت راست هدر: تاریخ روز به‌صورت خودکار
\rhead{\today}

% پایین صفحه: شماره صفحه
\cfoot{\thepage}

%---------------------------------------------------
% شروع سند
\begin{document}
	
	% صفحه عنوان (بدون هدر و فوتر)
	\thispagestyle{empty}
	
	\begin{center}
		
		% لوگوی دانشگاه
		\includegraphics[width=0.5\textwidth]{../../../images/image-E_2/image-1.png} \\[10pt]
		
		% عنوان سند
		\textbf{\LARGE عنوان: تمرین سری دو}\\[20pt]
		
		% نیم‌سال تحصیلی
		\textbf{\LARGE نیم‌سال تحصیلی: 4051}\\[10pt]
		
		% نام مدرس درس
		\textbf{\Large مدرس: دکتر مقدسین}\\[10pt]
		
		% مبحث تمرین
		\textbf{\Large مبحث تمرین: ضرب داخلی، ضرب برداری و ضرب سه‌گانه}\\[10pt]
		
		% مهلت تحویل تمرین
		\textbf{\Large مهلت تحویل: 24 مهر}
		
	\end{center}
	
	% صفحه جدید برای ادامه سند
	\newpage
	
	% ایجاد فهرست مطالب
	\tableofcontents
	
	\newpage
	
	%===================================================
	% سؤال اول
	%===================================================
\section{ضرب برداری در صفحه}

اگر
\[
\mathbf{P}
=
\hat{\mathbf e}_x P_x
+
\hat{\mathbf e}_y P_y
\]
و
\[
\mathbf{Q}
=
\hat{\mathbf e}_x Q_x
+
\hat{\mathbf e}_y Q_y
\]
دو بردار غیرموازی و غیرپادموازی در صفحه‌ی $xy$ باشند، نشان دهید که
\[
\mathbf{P}\times\mathbf{Q}
\]
در راستای محور $z$ قرار دارد.


%===================================================
% سؤال دوم
%===================================================

\section{ضرب داخلی و ضرب برداری}

ثابت کنید:
\[
(\mathbf{A}\times\mathbf{B})
\cdot
(\mathbf{A}\times\mathbf{B})
=
(AB)^2
-
(\mathbf{A}\cdot\mathbf{B})^2.
\]


%===================================================
% سؤال سوم
%===================================================

\section{اثبات اتحادهای مثلثاتی با استفاده از بردارها}

بردارهای زیر را در نظر بگیرید:

\[
\mathbf{P}
=
\hat{\mathbf e}_x\cos\theta
+
\hat{\mathbf e}_y\sin\theta,
\]

\[
\mathbf{Q}
=
\hat{\mathbf e}_x\cos\phi
-
\hat{\mathbf e}_y\sin\phi,
\]

و
\[
\mathbf{R}
=
\hat{\mathbf e}_x\cos\phi
+
\hat{\mathbf e}_y\sin\phi.
\]

با استفاده از بردارهای فوق، اتحادهای مثلثاتی زیر را اثبات کنید:

\[
\sin(\theta+\phi)
=
\sin\theta\cos\phi
+
\cos\theta\sin\phi,
\]

و

\[
\cos(\theta+\phi)
=
\cos\theta\cos\phi
-
\sin\theta\sin\phi.
\]


%===================================================
% سؤال چهارم
%===================================================

\section{بردار عمود بر دو بردار}

\subsection*{الف)}

برداری $\mathbf{A}$ بیابید که بر هر دو بردار

\[
\mathbf{U}
=
2\hat{\mathbf e}_x
+
\hat{\mathbf e}_y
-
\hat{\mathbf e}_z
\]

و

\[
\mathbf{V}
=
\hat{\mathbf e}_x
-
\hat{\mathbf e}_y
+
\hat{\mathbf e}_z
\]

عمود باشد.

\subsection*{ب)}

اگر علاوه بر شرط بالا بخواهیم $\mathbf{A}$ دارای اندازه‌ی واحد باشد، بردار $\mathbf{A}$ را به دست آورید.


%===================================================
% سؤال پنجم
%===================================================

\section{ضرب برداری بردارهای هم‌صفحه}

اگر چهار بردار $\mathbf{a}$، $\mathbf{b}$، $\mathbf{c}$ و $\mathbf{d}$ همگی در یک صفحه قرار داشته باشند، نشان دهید:

\[
(\mathbf{a}\times\mathbf{b})
\times
(\mathbf{c}\times\mathbf{d})
=
0.
\]

\textbf{راهنما:}

به راستای بردارهای حاصل از ضرب‌های برداری
\[
\mathbf{a}\times\mathbf{b}
\qquad\text{و}\qquad
\mathbf{c}\times\mathbf{d}
\]
توجه کنید.


%===================================================
% سؤال ششم
%===================================================

\section{محاسبه‌ی میدان مغناطیسی از نیروی لورنتس}

القای مغناطیسی $\mathbf{B}$ به وسیله‌ی معادله‌ی نیروی لورنتس تعریف می‌شود:

\[
\mathbf{F}
=
q(\mathbf{v}\times\mathbf{B}).
\]

با انجام سه آزمایش جداگانه، نتایج زیر به دست آمده است:

\[
\mathbf{v}
=
\hat{\mathbf e}_x,
\qquad
\frac{\mathbf{F}}{q}
=
2\hat{\mathbf e}_z
-
4\hat{\mathbf e}_y,
\]

\[
\mathbf{v}
=
\hat{\mathbf e}_y,
\qquad
\frac{\mathbf{F}}{q}
=
4\hat{\mathbf e}_x
-
\hat{\mathbf e}_z,
\]

و

\[
\mathbf{v}
=
\hat{\mathbf e}_z,
\qquad
\frac{\mathbf{F}}{q}
=
\hat{\mathbf e}_y
-
2\hat{\mathbf e}_x.
\]

با استفاده از نتایج این سه آزمایش، بردار القای مغناطیسی $\mathbf{B}$ را محاسبه کنید.


%===================================================
% سؤال هفتم
%===================================================

\section{اثبات اتحاد ژاکوبی برای ضرب برداری}

اتحاد ژاکوبی برای ضرب برداری را اثبات کنید:

\[
\mathbf{a}\times(\mathbf{b}\times\mathbf{c})
+
\mathbf{b}\times(\mathbf{c}\times\mathbf{a})
+
\mathbf{c}\times(\mathbf{a}\times\mathbf{b})
=
0.
\]


%===================================================
% سؤال هشتم
%===================================================

\section{نیروی مغناطیسی بین دو بار متحرک}

یک بار الکتریکی $q_1$ که با سرعت $\mathbf{v}_1$ حرکت می‌کند، طبق قانون بیو ـ ساوار القای مغناطیسی

\[
\mathbf{B}
=
\frac{\mu_0}{4\pi}
\frac{q_1}{r^2}
\left(
\mathbf{v}_1\times\hat{\mathbf r}
\right)
\]

را ایجاد می‌کند، که در آن $\hat{\mathbf r}$ یک بردار یکه در راستای خطی است که از $q_1$ به نقطه‌ای که میدان در آن اندازه‌گیری می‌شود، رسم شده است.

\subsection*{الف)}

نشان دهید نیروی مغناطیسی وارد بر بار دوم $q_2$ که با سرعت $\mathbf{v}_2$ حرکت می‌کند، به صورت ضرب سه‌گانه‌ی برداری زیر نوشته می‌شود:

\[
\mathbf{F}_2
=
\frac{\mu_0}{4\pi}
\frac{q_1q_2}{r^2}
\mathbf{v}_2
\times
\left(
\mathbf{v}_1\times\hat{\mathbf r}
\right).
\]

\subsection*{ب)}

نیروی مغناطیسی متناظر، یعنی نیروی $\mathbf{F}_1$ که بار $q_2$ بر بار $q_1$ وارد می‌کند، بنویسید.

بردار یکه‌ی شعاعی را تعریف کنید و سپس رابطه‌ی میان $\mathbf{F}_1$ و $\mathbf{F}_2$ را بررسی کنید.

\subsection*{پ)}

برای حالت خاصی که $q_1$ و $q_2$ در امتداد دو مسیر موازی و در کنار یکدیگر حرکت می‌کنند، $\mathbf{F}_1$ و $\mathbf{F}_2$ را محاسبه کنید.


%===================================================
% سؤال نهم
%===================================================

\section{ضرب داخلی دو ضرب برداری}

نشان دهید:

\[
\boxed{
(\mathbf{A}\times\mathbf{B})
\cdot
(\mathbf{C}\times\mathbf{D})
=
(\mathbf{A}\cdot\mathbf{C})
(\mathbf{B}\cdot\mathbf{D})
-
(\mathbf{A}\cdot\mathbf{D})
(\mathbf{B}\cdot\mathbf{C})
}.
\]


%===================================================
% سؤال دهم
%===================================================

\section{ضرب سه‌گانه‌ی برداری}

نشان دهید:

\[
\boxed{
(\mathbf{A}\times\mathbf{B})
\times
(\mathbf{C}\times\mathbf{D})
=
(\mathbf{A}\cdot(\mathbf{B}\times\mathbf{D}))\mathbf{C}
-
(\mathbf{A}\cdot(\mathbf{B}\times\mathbf{C}))\mathbf{D}
}.
\]


%===================================================
% سؤال امتیازی
%===================================================

\section{سؤال امتیازی: نمادگذاری لوی ـ چیویتا}

با استفاده از نمادگذاری لوی ـ چیویتا
\[
\varepsilon_{ijk}
\]
و قرارداد جمع اینشتین، اتحادهای برداری زیر را اثبات کنید.

\subsection*{الف)}

اثبات کنید:

\[
\boxed{
(\mathbf{A}\times\mathbf{B})
\cdot
(\mathbf{C}\times\mathbf{D})
=
(\mathbf{A}\cdot\mathbf{C})
(\mathbf{B}\cdot\mathbf{D})
-
(\mathbf{A}\cdot\mathbf{D})
(\mathbf{B}\cdot\mathbf{C})
}.
\]

\subsection*{ب)}

با استفاده از نمادگذاری لوی ـ چیویتا اثبات کنید:

\[
\boxed{
(\mathbf{A}\times\mathbf{B})
\times
(\mathbf{C}\times\mathbf{D})
=
(\mathbf{A}\cdot(\mathbf{B}\times\mathbf{D}))\mathbf{C}
-
(\mathbf{A}\cdot(\mathbf{B}\times\mathbf{C}))\mathbf{D}
}.
\]

در این سؤال، تمام مراحل تبدیل به نمایش اندیسی و ساده‌سازی عبارت‌ها را به‌طور کامل نشان دهید.


\vspace{20pt}

\begin{center}
	\textbf{\Large موفق باشید}
\end{center}

\end{document}